%%\documentstyle[twocolumn,jsaiac]{article}
%%\documentstyle[twocolumn,jsaiac]{article}
\documentclass[a4paper,twocolumn]{article}
\usepackage{jsaiac}
\usepackage{color}
\usepackage{comment}
\usepackage{url}
\usepackage[dvipdfmx]{xcolor}
\usepackage[dvips]{graphicx}

%%Title: 
\title{\LARGE \bf Deep Predictive Learning: Motion Learning Concept\\
\LARGE \bf inspired by Cognitive Robotics}

%%Contact address: 
\address{Tetsuya Ogata, Waseda University, 3-4-1 Ookubo, Shinjuku-ku, Tokyo, 169-8555, Japan, {\tt\small ogata@waseda.jp}\\
\hspace{-2.5mm}${\dagger}$ All authors have contributed equally to this work.
}

%%Name of authors:
\author{%
Kanata Suzuki$^{\dagger}$\first\second
\and
\hspace{6mm}
\and
Hiroshi Ito$^{\dagger}$\first\third
\and
\hspace{6mm}
\and
Tatsuro Yamada$^{\dagger}$\first\fourth
\and
\hspace{6mm}
\and
Kei Kase$^{\dagger}$\first
\and
\hspace{6mm}
\and
Tetsuya Ogata$^{\dagger}$\first
}

%%Affiliations:
\affiliate{
\first{} Waseda University
\and
\second{} Fujitsu Limited
\and
\third{} Hitachi, Ltd.
\and
\fourth{} Panasonic Corporation
}

%%
%\Vol{28}        %% <-- 25
%\session{0A0-00}%% <-- ID

\begin{abstract}
%実世界でロボットを稼働させるためには，動作モデルと現実のギャップを限られたデータから埋めることが重要である．
%深層学習を用いた学習型アプローチは環境認識と動作生成をEnd-to-Endに学習することで，特徴量設計コストを削減しながら，多様な状況に汎化していくことが期待されている．
%しかしながら，学習データの収集コストが課題であり，物理的な接触を伴うロボットの試行錯誤には時間と人的リソースが必要不可欠である．
%上記の課題を解決するために我々は，予測モデルの不完全性を前提とし，ロボットの感覚運動ダイナミクスを一体となって予測する動作学習コンセプトである「深層予測学習」を提唱する．
%深層予測学習は，生体が実世界と脳の予測誤差が最小となるように振る舞うことを説明する予測符号化理論にinspireされている．
%ロボットは感覚運動情報に基づき近未来の状況を予測し，現実との誤差を最小とするようにモデルを学習する．
%結果として，ロボットは学習時と現実の差を許容しながらリアルタイムに動作を調整し，未学習の状況下でも柔軟に作業を実行することが可能である．
%また，モデルに埋め込まれた動作ダイナミクスを組み合わせながら，幅広いタスク動作を実現することが期待される．
%本論文では，深層予測学習の概念，実装方法，更に実ロボットへの適用事例について述べる．
%To deploy robots in the real world, it is crucial to bridge the gap between the motion model and reality with limited data. Deep learning is expected to generalize to diverse situations while reducing feature design costs through end-to-end learning for environment recognition and motion generation. However, data collection for model training is highly costly, and time and human resources are essential for robot trial-and-error with physical contact. To solve the above problems, we propose "Deep Predictive Learning", a motion learning concept that predicts the robot's sensorimotor dynamics as a whole, assuming imperfections in the prediction model. This concept is inspired by predictive coding theory. It is based on the fundamental strategy of predicting the near-future sensorimotor states of robots and online minimization of the prediction error between the real world and the model. Based on the acquired sensor information, the robot can adjust its behavior in real time, tolerating the difference between the learning experience and reality. In addition, the robot is expected to perform a wide range of tasks while combining the motion dynamics embedded in the model. This paper describes the proposed concept, how it is implemented, and examples of its application to real robots.
Bridging the gap between motion models and reality is crucial by using limited data to deploy robots in the real world. Deep learning is expected to be generalized to diverse situations while reducing feature design costs through end-to-end learning for environmental recognition and motion generation. However, data collection for model training is costly, and time and human resources are essential for robot trial-and-error with physical contact. We propose “Deep Predictive Learning,” a motion learning concept that predicts the robot’s sensorimotor dynamics, assuming imperfections in the prediction model. The predictive coding theory inspires this concept to solve the above problems. It is based on the fundamental strategy of predicting the near-future sensorimotor states of robots and online minimization of the prediction error between the real world and the model. Based on the acquired sensor information, the robot can adjust its behavior in real time, thereby tolerating the difference between the learning experience and reality. Additionally, the robot was expected to perform a wide range of tasks by combining the motion dynamics embedded in the model. This paper describes the proposed concept, its implementation, and examples of its applications in real robots.
The code and documents are available at: \url{https://ogata-lab.github.io/eipl}
\end{abstract}


%\setcounter{page}{1}
\def\Style{``jsaiac.sty''}
\def\BibTeX{{\rm B\kern-.05em{\sc i\kern-.025em b}\kern-.08em%
 T\kern-.1667em\lower.7ex\hbox{E}\kern-.125emX}}
\def\JBibTeX{\leavevmode\lower .6ex\hbox{J}\kern-0.15em\BibTeX}
\def\LaTeXe{\LaTeX\kern.15em2$_{\textstyle\varepsilon}$}

\begin{document}
\maketitle



\section{Introduction}

%"The Moravec Paradox" is currently one of the significant challenges remaining in artificial intelligence technology \cite{Moravec1988mind}. This paradox refers to the contradiction where tasks that even children can perform semi-unconsciously are exceedingly difficult for the latest artificial intelligence and robots compared to complex logical reasoning, such as puzzles or chess. Tasks involving independent movement like walking or running made significant technological advancements in the 2000s \cite{Garcia1996,Hirai1998}, while perception tasks such as vision and voice recognition saw substantial progress in the 2010s through deep learning techniques~\cite{Krizhevsky2012nips,He2016cvpr,Ren2016pami,Redmon2016cvpr,Hinton2013spm,Graves2014icml}. However, performing a variety of tasks using a single common hand, just like humans, remains extremely challenging even with the latest technology available today.
"The Moravec Paradox" is currently one of the significant challenges in artificial intelligence technology \cite{Moravec1988mind}. This paradox refers to the contradiction in which tasks that even children can perform semi-unconsciously are exceedingly difficult for the latest artificial intelligence and robots compared to complex logical reasoning tasks, such as puzzles or chess. Tasks involving independent movements such as walking or running made significant technological advancements in the 2000s \cite{Garcia1996,Hirai1998}. In contrast, perception tasks such as vision and voice recognition saw substantial progress in the 2010s through deep learning techniques~\cite{Krizhevsky2012nips,He2016cvpr,Ren2016pami,Redmon2016cvpr,Hinton2013spm,Graves2014icml}. However, performing a variety of tasks using a single common hand, just like humans, remains extremely challenging even with the latest technology available today.



\subsection{Experience-based Robotics}

%The use of deep learning methods is expanding in the fields of both environmental recognition and motion generation.  Examples of this include extracting graspable regions of objects~\cite{Pinto2016icra,Mahler2017dex,Mousavian2019iccv} and combinatorial exploration of action patterns through reinforcement learning~\cite{Levine2016jmlr,Levine2017ijrr}. These research examples represent an inductive approach to directly acquiring models of the environment and trajectory generation from data. In the context of robot learning, data can be generated or acquired as a result of its bottom-up exploratory actions. Sensory-motor information obtained from active interaction, or "experience," plays an important role in robot learning. For example, while conventional learning data is annotated with labels provided by humans, the experience can be utilized in the robot's own motor skill learning even in the absence of teacher labels. Another important perspective that emerges in this concept called experience-based robotics is the functional separation of environmental perception and motion generation. Considering the now common end-to-end approach in processing images, audio, and natural language using deep learning, it becomes evident that the clear separation of these two functions could also be a subject of learning.
Deep learning methods are expanding in the environmental recognition and motion generation fields. Examples include the extraction of graspable regions of objects~\cite{Pinto2016icra,Mahler2017dex,Mousavian2019iccv} and the combinatorial exploration of action patterns through reinforcement learning~\cite{Levine2016jmlr,Levine2017ijrr}. These research examples represent an inductive approach for directly acquiring models of the environment and trajectory generation from data. In robot learning, data can be generated or acquired through bottom-up exploratory actions. Sensory-motor information obtained from active interaction, or "experience," plays an important role in robot learning. For example, while conventional learning data are annotated with labels provided by humans, experience can be utilized in a robot's motor-skill learning, even in the absence of teacher labels. Another important perspective from this concept, called experience-based robotics, is the functional separation between environmental perception and motion generation. Considering the common end-to-end approach for processing images, audio, and natural language using deep learning, it is evident that a clear separation of these two functions could also be a subject of learning.


%The concept of recognizing and generating partial "functions" of autonomous systems here is important, as it is believed to emerge gradually from the interaction of the body, environment, neural circuits, and the continuous dynamics of end-to-end learning, as proposed in cognitive developmental robotics~\cite{Asada2015} or brain-inspired learning~\cite{Haruno2001}. Particularly, to deploy robots in the real world, it is crucial to bridge the gap between the motion model and reality with limited data. While deep learning-based approaches show improvements in task accuracy through large amounts of data and trials, data collection for model training is highly costly. In particular, robot motion data depends on hardware, making it uncommon to have access to extensive data for model tuning. Although research is conducted to transfer learning results from simulation environments to the real world, modeling nonlinear physical phenomena remains challenging and is not entirely solvable. It is necessary to achieve a wide range of task behaviors by creatively combining representations derived from the limited motion experiences of robots.
The concept of recognizing and generating partial "functions" of autonomous systems here is important, as it is believed to emerge gradually from the interaction of the body, environment, neural circuits, and the continuous dynamics of end-to-end learning, as proposed in cognitive developmental robotics~\cite{Asada2015} or brain-inspired learning~\cite{Haruno2001}. In particular, to deploy robots in the real world, it is crucial to bridge the gap between motion models and reality using limited data. Although deep learning-based approaches show improvements in task accuracy through large amounts of data and trials, data collection for model training is costly. In particular, robot motion data depend on the hardware, making accessing extensive data for model tuning uncommon. Although research has been conducted to transfer learning results from simulation environments to the real world, modeling nonlinear physical phenomena remains challenging and unsolvable. It is necessary to achieve a wide range of task behaviors by creatively combining representations derived from the limited motion experience of robots.



\subsection{Proposed Robot Learning Concept}

%In this paper, we discuss the concept of "deep predictive learning," a methodology for applying deep learning techniques to robotics, along with its implementation methods. This concept, inspired by predictive coding theory~\cite{Rao1999, Friston2006jpp, Friston2010nrn}, is based on the fundamental strategy of predicting the near-future sensory-motor states of robots and online minimization of the prediction error between the real world and the model. Furthermore, by embedding the temporal relationships between sensory-motor information in systems where the environment and body interact dynamically, the system transitions to other motion dynamics based on prediction error during generation, enabling the execution of high-dimensional, long-term tasks.
This study discusses the concept of "deep predictive learning," a methodology for applying deep learning techniques to robotics, along with its implementation methods. This concept, inspired by predictive coding theory~\cite{Rao1999, Friston2006jpp, Friston2010nrn}, is based on the fundamental strategy of predicting the near-future sensory-motor states of robots and online minimization of the prediction error between the real world and the model. Furthermore, by embedding the temporal relationships between sensory and motor information in systems where the environment and body interact dynamically, the system transitions to other motion dynamics based on prediction errors during generation, thereby enabling the execution of high-dimensional, long-term tasks.


%In Section 2 of this paper, we discuss previous research and its challenges related to the application of deep learning in robotics and describe the positioning of deep predictive learning. Section 3 covers the background of deep predictive learning, including predictive coding and the free energy principle, and outlines strategies for applying them to robot motion learning. We also discuss specific implementation approaches based on deep learning techniques and examine them from the perspective of cognitive robotics, including dynamic cognitive systems and action-normative systems. In Section 4, we summarize the published studies that we have developed using the DPL concept from two perspectives: scalability by combining multiple motions, and application to complex manipulation tasks. Finally, in Section 5, we provide a summary and discuss future prospects.
Section 2 discusses previous research and the challenges related to applying deep learning in robotics and describes the positioning of deep predictive learning. Section 3 covers the background of deep predictive learning, including predictive coding and the free-energy principle, and outlines strategies for applying them to robot motion learning. We also discuss specific implementation approaches based on deep learning techniques and examine them from the perspective of cognitive robotics, including dynamic cognitive and action-normative systems. Section 4 summarizes the published studies we developed using the proposed concept from two perspectives: scalability by combining multiple motions and application to complex manipulation tasks. Finally, in Section 5, we summarize and discuss prospects.




\section{Deep Learning in Robotics Fields}

%In this Section, we will review previous research on the application of deep learning in robotics and establish the context for our proposed concept of deep predictive learning.
In this section, we review previous research on the application of deep learning in robotics and establish the context for our proposed deep predictive learning concept.



\subsection{Recognition}

%The term "deep learning" was first introduced in the field of robotics at the IEEE/RSJ IROS 2017 conference, which was relatively late compared to other domains. Many early applications of deep learning in robotics focused on areas such as object detection in robot vision \cite{Ren2016pami,Redmon2016cvpr} and research in conversational speech recognition \cite{Hinton2013spm,Graves2014icml}. In traditional robotic intelligence design, the process typically involves the following steps: (1) recognizing the environment or objects, (2) correlating recognition results with pre-provided 2D maps or 3D CAD models, (3) understanding the current positions, object positions, and orientations, and (4) planning actions based on the situation. Deep learning significantly enhanced the first step, which is recognition capability. By learning labels like graspable regions that were previously challenging to model, it became possible to recognize a wide variety of general objects.
The term "deep learning" was introduced in robotics at the IEEE/RSJ IROS 2017 conference, which was relatively late compared with other domains. Several early applications of deep learning in robotics focused on object detection in robot vision~\cite{Ren2016pami,Redmon2016cvpr} and conversational speech recognition~\cite{Hinton2013spm,Graves2014icml}. In traditional robotic intelligence design, the process typically involves the following steps: (1) recognizing the environment or objects, (2) correlating the recognition results with pre-provided 2D maps or 3D CAD models, (3) understanding the current positions, object positions, and orientations, and (4) planning actions based on the situation. Deep learning significantly enhances the recognition capability in the first step. Learning labels, such as graspable regions that were previously challenging to model, made recognizing a wide variety of general objects possible.


%However, the burden of manual labeling by humans is substantial, and relying solely on image-based learning does not guarantee successful object grasping. Grasping performance depends on a combination of visual information, object attributes, and the robot's physical characteristics. Factors like the gripper's force, friction conditions, and object deformations must all be considered, and there are countless such elements to account for.
However, the burden of manual labeling by humans is substantial, and relying solely on image-based learning does not guarantee successful object grasping. The grasping performance depends on a combination of visual information, object attributes, and physical characteristics of the robot. Factors such as the gripper force, friction conditions, and object deformations must be considered, as well as countless elements.



\subsection{Robot Motion Learning}

%In response to the developments in perception systems mentioned above, there has been active research~\cite{Noda2014multimodal,Lee2019icra,lin2022tactilegym2} in End-to-End methods for predicting robot motion trajectories directly from multi-modal input sensor data (such as visual, auditory, tactile, etc.). Within the realm of learning-based approaches for robot motion, deep learning-based methods can be broadly categorized into two groups.
In response to the developments mentioned above in perception systems, active research has been conducted ~\cite{Noda2014multimodal,Lee2019icra,lin2022tactilegym2} on end-to-end methods for predicting robot motion trajectories directly from multimodal input sensor data (such as visual, auditory, tactile, etc.). Deep learning-based methods can be broadly categorized into two groups within learning-based approaches to robot motion.



\subsubsection{Reinforcement Learning}

%In deep reinforcement learning, the agent predicts rewards for combinations of high-dimensional inputs related to the environment and several executable action outputs by the robot. With each attempt by the agent, rewards are provided from the environment, and the model updates its parameters to acquire a policy that maximizes rewards~\cite{Levine2016jmlr,Levine2017ijrr,Haarnoja2017icml}. The application of deep reinforcement learning in games has been extremely successful~\cite{Mnih2013playing,Hessel2018rainbow}, and its application to robotics is a natural progression. Its greatest appeal lies in the significant reduction of the traditional process of designing models for the environment or objects, which designers had to provide. This has led to the proposal of various task-specific methods~\cite{Tsurumine2019ras,Mahler2019learning}. Levine and colleagues replaced previously required image feature engineering with deep learning and built an end-to-end model from images to actions~\cite{Levine2017ijrr}. On the other hand, in~\cite{lee2023pi,rlscale2023arxiv}, efforts were made to enhance the performance of general agents by training multiple tasks with a single model. However, achieving the optimal model for online learning in a real environment remains a challenge, as it requires a large number of trials and resources. There is also the approach of collecting a large amount of data with human resources~\cite{dasari2019robonet,rt12022arxiv}, but this approach has limitations in terms of data diversity.
In deep reinforcement learning, the agent predicts rewards for combinations of high-dimensional inputs related to the environment and several executable action outputs from the robot. With each attempt by the agent, the environment provides rewards, and the model updates its parameters to acquire a policy that maximizes the rewards~\cite{Levine2016jmlr,Levine2017ijrr,Haarnoja2017icml}. The application of deep reinforcement learning in games has been extremely successful~\cite{Mnih2013playing,Hessel2018rainbow}, and its application in robotics is a natural progression. Its greatest appeal lies in significantly reducing the traditional process of designing models for the environment or objects that designers must provide. This innovation has led to various task-specific methods~\cite{Tsurumine2019ras,Mahler2019learning}. Levine et al. replaced the previously required image feature engineering with deep learning and built an end-to-end model from images to actions~\cite{Levine2017ijrr}. However, in~\cite{lee2023pi,rlscale2023arxiv}, efforts were made to enhance the performance of general agents by training multiple tasks with a single model. However, achieving an optimal model for online learning in a real environment remains a challenge, as it requires many trials and resources. There is also an approach for collecting a large amount of data using human resources~\cite{dasari2019robonet,rt12022arxiv}, but this approach has limitations regarding data diversity.


%As a solution to the above issues, there is the Sim2Real approach that uses simulation~\cite{James2019cvpr,Chebotar2019icra}. By making simulation closer to reality~\cite{allevato2020tunenet} and transforming it~\cite{rao2020rl}, it supplements the lack of data. Domain randomization, especially regarding visual and physical parameters~\cite{Tobin2017}, is utilized in tasks such as multi-fingered hand manipulation~\cite{openai2018learning} and rough terrain walking of quadruped robots~\cite{lee2020learning}. In~\cite{lee2020learning}, teacher-student learning using environment information obtained only in simulation has been successful in real-world motion control, which is quite intriguing. Recent advancements in physical simulation technology can replicate scenes equivalent to the real world, but detailed modeling of contact states remains a challenge~\cite{drake,Narang2022rss}. Currently, this problem is being discussed as one of the subjects in the field of Surface Science.
The Sim2Real approach uses simulations to address these issues~\cite{James2019cvpr,Chebotar2019icra}. Making simulations closer to reality~\cite{allevato2020tunenet} and transforming them~\cite{rao2020rl} can supplement the lack of data. Domain randomization, especially regarding visual and physical parameters~\cite{Tobin2017}, is utilized in tasks such as multi-fingered hand manipulation~\cite{openai2018learning} and rough terrain walking of quadruped robots~\cite{lee2020learning}. In~\cite{lee2020learning}, teacher-student learning using environmental information obtained only in simulations was successful in real-world motion control, which is quite intriguing. Recent advancements in physical simulation technology can replicate scenes equivalent to the real world; however, detailed modeling of contact states remains a challenge~\cite{drake,Narang2022rss}. This problem is currently being discussed as a subject in Surface Science.


%Moreover, traditional reinforcement learning typically requires the pre-designing of reward functions, which poses challenges as the model's performance is heavily dependent on it. Inverse reinforcement learning introduces a process of estimating reward functions from human demonstrations~\cite{Peng2018,Tsurumine2019humanoids,Wang2021}, reducing the burden of manually designing rewards based on task domain knowledge. However, aspects such as policy behavior design still remain, and the flexibility of applied tasks on real machines is limited.
Moreover, traditional reinforcement learning typically requires pre-designing reward functions, which poses challenges because the model's performance depends heavily on it. Inverse reinforcement learning introduces a process for estimating the reward functions from human demonstrations~\cite{Peng2018,Tsurumine2019humanoids,Wang2021}, thereby reducing the burden of manually designing rewards based on task domain knowledge. However, aspects such as policy behavior design remain, and the flexibility of the tasks applied to real machines is limited.



\subsubsection{Behavioral Cloning}

%In recent years, there has been growing interest in offline learning to address the challenges of data scale arising from online policy updates in conventional reinforcement learning. Particularly, in Behavioral Cloning (or Imitation Learning), learning is conducted to replicate motion sequence data (joint angles, velocities, torques, etc.) obtained from robots~\cite{yu2018one,Zhao2023rss,jang2021bc,Lynch2022,guhur2023instruction,chi2023diffusionpolicy}. Zhao et al. used a low-cost, fine-grained teleoperation system and employed Transformers to enable learning across various tasks~\cite{Zhao2023rss,fu2024mobile}. In this framework, it is possible to directly acquire a motion generation model from motion experience if effective data for task execution is available. Additionally, the feature of not requiring complex reward design is highlighted. As a result, it has been applied to various high-difficulty tasks such as flexible object manipulation~\cite{Kawaharazuka2019icra} and human-cooperative tasks~\cite{Sasagawa2020ral}. 
Recently, there has been growing interest in offline learning to address the challenges of data scale arising from online policy updates in conventional reinforcement learning. In Behavioral Cloning (or Imitation Learning), learning is conducted to replicate the motion sequence data (joint angles, velocities, torques, etc.) obtained from robots~\cite{yu2018one,Zhao2023rss,jang2021bc,Lynch2022,guhur2023instruction,chi2023diffusionpolicy}. Zhao et al. used a low-cost fine-grained teleoperation system and employed transformers to enable learning across various tasks~\cite{Zhao2023rss,fu2024mobile}. In this framework, it is possible to acquire a motion generation model directly from the motion experience if effective data for task execution are available. In addition, the feature of not requiring a complex reward design is highlighted. Consequently, it has been applied to various high-difficulty tasks, such as flexible object manipulation~\cite{Kawaharazuka2019icra} and human-cooperative tasks~\cite{Sasagawa2020ral}. 


%As related technologies, we should also mention offline reinforcement learning. Offline reinforcement learning aims to acquire a good policy solely from a pre-collected dataset, and these studies can be interpreted as reward-conditioned Behavioral Cloning~\cite{actionablemodels2021arxiv,ma2022vip}. These approaches seek to improve learning efficiency by utilizing past experience data~\cite{kalashnikov2018qt} or by simultaneously collecting data from multiple robots~\cite{kalashnikov2021mt}. In particular, Chen et al. proposed a sequentially predictable offline reinforcement learning model using Transformers~\cite{chen2021decision}. They predict actions from a sequence of tokens representing rewards, states, and actions, achieving high performance in long-term tasks using the expected cumulative future rewards.
Regarding related technologies, we should also mention offline reinforcement learning. Offline reinforcement learning aims to acquire a good policy solely from a pre-collected dataset, and these studies can be interpreted as reward-conditioned Behavioral Cloning~\cite{actionablemodels2021arxiv,ma2022vip}. These approaches seek to improve the learning efficiency by utilizing experience data~\cite{kalashnikov2018qt} or by simultaneously collecting data from multiple robots~\cite{kalashnikov2021mt}. In particular, Chen et al. proposed a sequentially predictable offline reinforcement learning model using transformers~\cite{chen2021decision}. They predicted actions from a sequence of tokens representing rewards, states, and actions and achieved high performance in long-term tasks using expected cumulative future rewards.


%Many of the studies mentioned so far focus on modeling robot action sequences alone. In contrast, the "Predictive Learning" approach aims to perform self-supervised learning, including the external environment~\cite{lotter2016deep}. Lotter et al. proposed PredNet, which mimics predictive coding processes in the cerebral cortex. PredNet learns to minimize the error between predicted and actual values in video prediction tasks. These are implementations of the predictive coding theory~\cite{Rao1999} using deep learning models but do not consider interaction with the external environment in robotics.
Many of the studies mentioned thus far have focused solely on modeling robot action sequences. In contrast, the "Predictive Learning" approach aims to perform self-supervised learning, including the external environment~\cite{lotter2016deep}. Lotter et al. proposed PredNet, which mimics the predictive coding processes in the cerebral cortex. PredNet learns to minimize the error between the predicted and actual values in video prediction tasks. These are implementations of predictive coding theory~\cite{Rao1999} using deep learning models, but they do not consider interaction with the external environment in robotics.


%Our proposed concept, from the perspective of cognitive robotics~\cite{Asada2015}, incorporates predictive coding in sensorimotor dynamics. The crucial idea in our proposal is that the learning target of autonomous systems should not be a "sensor-to-motion map" but rather a "continuous loop of integrated sensorimotor dynamics." In systems where the environment and the body dynamically interact, the distinction between sensation and motion is always ambiguous. What exists there is a continuous flow where the current something leads to the future outcome. It is crucial to have an orientation to fill the error of the predictive flow that never completely fills the self and external. This kind of predictive learning system (forward model) and the corresponding motion generation system (inverse model) have been emphasized for a long time, but practical robots explicitly designed to work with this concept in the real world are still limited.
Our proposed concept, from the perspective of cognitive robotics~\cite{Asada2015}, incorporates predictive coding into sensorimotor dynamics. The crucial idea in our proposal is that the learning target of autonomous systems should not be a "sensor-to-motion map" but rather a "continuous loop of integrated sensorimotor dynamics." The distinction between sensation and motion is always ambiguous in systems where the environment and body interact dynamically. A continuous flow exists in which the current leads to future outcomes. It is crucial to have an orientation that fills the error of the predictive flow, which never completely fills the self and the external. This type of predictive learning system (forward model) and the corresponding motion generation system (inverse model) have been emphasized for a long time. However, practical robots designed to work with this concept in the real world remain limited.






%-------------------- Fig --------------------
\begin{figure}[t]
\centering
\includegraphics[width=1.0\columnwidth]{fig_dpl.eps}
\caption{
Overview of the free energy principle~\cite{Friston2010nrn}.
}
\label{fig_dpl}
\hspace{10mm}
\end{figure}



\section{Deep Predictive Learning}

%This section explains the four strategies that form the basis of the deep predictive learning (DPL) concept, along with the free energy principle and predictive coding. Then, as an implementation of the DPL concept, we describe the phases for building a basic predictive model and three particularly important elements in that phase. Also, we discuss the concept from the perspective of embodiment~\cite{Brooks1991intelligence} and dynamic cognitive system~\cite{Tani1996model}.
This section explains the four strategies that form the basis of the deep predictive learning (DPL) concept, free-energy principle, and predictive coding. Then, as an implementation of the DPL concept, we describe the phases for building a basic predictive model and the three particularly important elements in that phase. In addition, we discussed this concept from the perspectives of embodiment~\cite{Brooks1991intelligence} and dynamic cognitive systems~\cite{Tani1996model}.



\subsection{Main Concept and Policies}

%The Deep Neural Network (DNN) framework uses large amounts of data to optimize models, and the goal is for that model to function appropriately even with new data. However, it is usually not easy to obtain an optimal model because it is impossible to know which state will provide the greatest generalizability. Especially in robots, it can be assumed that they are constantly exposed to unexpected situations because they operate in the real environment. Therefore, an approach is required that assumes that the learning is incomplete and considers how to deal with prediction errors.
The Deep Neural Network (DNN) framework uses large amounts of data to optimize models, and the goal is for the model to function appropriately, even with new data. However, obtaining an optimal model is usually difficult because it is impossible to determine the state that provides the greatest generalizability. It can be assumed that robots are constantly exposed to unexpected situations because they operate in real environments. Therefore, an approach that assumes that learning is incomplete and considers how to manage prediction errors is required.


%Here, there is an information theory of the brain called the free energy principle~\cite{Friston2006jpp,Friston2010nrn} proposed by K. Friston that can help tackle the above problem. The free energy principle states that the perception, learning, and behavior of living things are determined by minimizing a cost function called variational free energy (Figure~\ref{fig_dpl}). We try to summarize it below.
There is an information theory of the brain called the free energy principle~\cite{Friston2006jpp,Friston2010nrn} proposed by K. Friston, which can help address the above problem. The free-energy principle states that living things' perception, learning, and behavior are determined by minimizing a cost function called variational free energy (Figure~\ref{fig_dpl}). We summarize this as follows:
\begin{equation}
    FreeEnergy = -\log p(x|z) + D_{KL}[q(z|x)||p(z)] \label{eq_free}
\end{equation}
%The above corresponds to the negative Evidence Lower Bound in general machine learning models, and this equation itself is not new. However, by considering the interpretation of the equation and the multiple parameters involved, we can discuss the design of a system with a concrete body. The brain is a predictor and has a model $p(z)$ as a priori belief about the perception $z$ from the environment that causes the sensory input $x$. In addition, perception $z$ includes not only the hidden states of exteroceptive sensations such as vision and touch but also the hidden states of proprioception (e.g., joints) and interoceptive sensations related to kinesthetic and visceral sensations. The brain generates a sensory prediction $g_w(z)$ based on this perception $z$, and obtains a posteriori belief $q(z|x)$ by receiving the current interoceptive sensory input $x$. After that, the brain continuously generates predictions of the next sensory input based on the updated posterior beliefs and updates the beliefs. The predicted sensory inputs include a prediction of accuracy (the reciprocal of the sensation variance $v_x$). In other words, the first term in the equation \ref{eq_free} refers to the prediction error weighted by this accuracy, and the second term refers to the influence of the complexity of the structure of the prior belief (model).
The above equation corresponds to the negative Evidence Lower Bound in general machine learning models, which is not new. However, by considering the interpretation of the equation and the multiple parameters involved, we can discuss the design of a system with a concrete body. The brain is a predictor with a model $p(z)$ as an a priori belief about the perception $z$ from the environment that causes sensory input $x$. In addition, perception $z$ includes not only the hidden states of exteroceptive sensations, such as vision and touch but also the hidden states of proprioception (e.g., joints) and interoceptive sensations related to kinesthetic and visceral sensations. The brain generates a sensory prediction $g_w(z)$ based on this perception $z$ and obtains a posteriori belief $q(z|x)$ by receiving the current interoceptive sensory input $x$. Subsequently, the brain continuously generates predictions for the next sensory input based on updated posterior beliefs and updates those beliefs. The predicted sensory inputs included the prediction accuracy ( reciprocal of the sensation variance $v_x$). In other words, the first term in equation \ref{eq_free} refers to the prediction error weighted by this accuracy. The second term refers to the influence of the complexity of the prior belief (model) structure.


%Strategies for reducing prediction errors are called "predictive coding," and can be divided into four main types. By combining these strategies, the DPL realizes online motion generation tailored to the environment.
Strategies for reducing prediction errors are called "predictive coding" and can be classified into four main types. The DPL realizes online motion generation that is tailored to the environment by combining these strategies.


\begin{description}
\item[Training]\mbox{}\\
%The basic strategy is to optimize the entire model. The model that generates sensory predictions is trained, and the internal weight parameter $w$ is changed according to the prediction error. However, training the DNN model is costly and difficult to implement under time constraints. This problem is particularly serious in deep learning approaches, as the model size tends to be large. Furthermore, depending on the real robot situation, frequent model modifications may be inappropriate.
The basic strategy is to optimize the entire model. The model that generates the sensory predictions is trained, and the internal weight parameter $w$ changes based on the prediction error. However, training a DNN model is expensive and difficult under time constraints. This problem is particularly serious in deep learning approaches because the model size tends to be large. Furthermore, frequent model modifications may be inappropriate depending on the robot's situation.


\item[Modification of Posterior Belief]\mbox{}\\
%The second strategy is to modify the hypothetical posterior belief $q(z)$. The values of sensory input are modified in a closed manner based on the predicted perception, and as a result, the prediction error of the model is reduced. For example, when humans perceive the environment, they often interpret the real world as an afterthought based on previous images in their brains. And, the brain is unaware of the errors. This strategy is a method of adapting to changes in the real world by using predictive biases to create a virtual perception that differs from reality (Section~\ref{sec_generation}).
The second strategy is to modify hypothetical posterior belief $q(z)$. The values of the sensory input were modified in a closed manner based on the predicted perception; as a result, the prediction error of the model was reduced. For example, when humans perceive an environment, they often interpret the real world as an afterthought based on previous brain images. Additionally, the brain is unaware of these errors. This strategy is a method of adapting to real-world changes using predictive biases to create a virtual perception that differs from reality (Section~\ref{sec_generation}).


\item[Active Inference]\mbox{}\\
%The third strategy is to use the body to influence the external environment and change the sensory input $x$ to reduce prediction errors. This operation is called active inference in the context of the free energy principle. There are methods such as manipulating the environment by oneself or changing the viewpoint by changing one's position. This strategy is essential for autonomous agents that have bodies in the real world, such as robots. We implement the behavioral transitions for performing posterior beliefs and active inference using attractor dynamics embedded in the internal state of the model (Sectoin~\ref{sec_imple}).
The third strategy is to use the body to influence the external environment and change the sensory input $x$ to reduce the prediction errors. This operation is called active inference in the context of the free energy principle. These methods include manipulating the environment by oneself or changing viewpoints by changing one's position. This strategy is essential for autonomous agents, such as robots, with bodies in the real world. We implemented behavioral transitions to perform posterior belief and active inference using attractor dynamics embedded in the model's internal state (Sectoin~\ref{sec_imple}).


\item[Modification of Prediction Accuracy]\mbox{}\\
%The fourth strategy is to reduce the apparent error by setting the prediction accuracy $1/v_x$ of sensory input $x$ small and multiplying it by the weight of the prediction error. This operation, which increases the precision of sensory input, is associated with directing attention. From the perspective of predictive encoding, this strategy can be interpreted as placing emphasis on prediction errors and increasing prediction corrections when sensory input is reliable (when attention is directed). On the other hand, it can be interpreted that when sensory input is unreliable (when attention is not directed), one ignores the prediction error and maintains one's inference. The implementation example corresponds to Section~\ref{sec_attention}.
The fourth strategy involves reducing the apparent error by setting the prediction accuracy $1/v_x$ of the sensory input $x$ and multiplying it by the weight of the prediction error. This operation, which increases the precision of sensory input, is associated with directing attention. From the perspective of predictive encoding, this strategy can be interpreted as emphasizing prediction errors and increasing prediction corrections when sensory input is reliable (when attention is directed). However, when sensory input is unreliable (when attention is not directed), the prediction error is ignored, and the inference is maintained. The implementation example corresponds to that presented in Section~\ref{sec_attention}.
\end{description}


%-------------------- Fig --------------------
\begin{figure*}[t]
\centering
\includegraphics[width=17cm]{1_overview.eps}
\caption{
%An example of a deep predictive learning framework, consisting of a feature extraction unit and a time-series prediction unit.
An example of a deep predictive learning framework consisting of a feature extraction unit and a time-series prediction unit.
}
\label{fig_overview}
\end{figure*}



\subsection{Implementation}

\subsubsection{Learning Sensorimotor Dynamics}
\label{sec_imple}

%We explain the construction phase of the DPL model to realize the above-mentioned strategies, divided into training data collection and model learning. Figure~\ref{fig_overview} shows the overall flow including the generation phase described in section \ref{sec_generation}.
We explain the construction phase of the DPL model to realize the strategies mentioned above, which are divided into training data collection and model learning. Figure~\ref{fig_overview} illustrates the overall flow, including the generation phase described in Section \ref{sec_generation}.


\textbf{Collecting Training Dataset:}
%In the DPL concept, a model to generate robot motions is acquired in a data-driven manner by directly using sensorimotor information from the robot (time-series data consisting of motion information, sensor information, etc.) as training data. Therefore, the training data must include information about the interaction between the robot's body and the environment. Furthermore, since the quantity and quality of training data greatly affect model performance, it is desirable to be able to efficiently collect high-quality demonstration data.
In the DPL concept, a model for generating robot motions is acquired in a data-driven manner by directly using sensorimotor information from the robot (time-series data consisting of motion and sensor information, etc.) as training data. Therefore, the training data must include information on the robot's body and environment interaction. Furthermore, because the quantity and quality of training data greatly affect model performance, it is desirable to collect high-quality demonstration data efficiently.


%In the first step, training data is collected by actually having the robot perform the desired task and recording motion information such as joint angles, and sensor information such as camera images at a fixed sampling rate. Typical teaching methods for robot motion include programming description, direct teaching, and teleoperation. Describing the motion in advance using a robot programming language is the simplest teaching method, but the description becomes complicated for long-horizon tasks. On the other hand, by having a human remotely control the robot and teaching it its motions, it is possible to obtain training data without precise modeling or parameter adjustment. Especially, the method of remotely controlling the manipulator from the actual robot's perspective (Wizard of Oz method) is desirable because it allows the robot to intuitively teach human operation skills for tasks. By controlling the robot as if they were controlling their own bodies, workers can interact naturally with the environment. The acquired teaching data are expected to include the information necessary for motion learning.
In the first step, training data are collected by having the robot perform the desired task and record motion information, such as joint angles, and sensor information, such as camera images, at a fixed sampling rate. Typical teaching methods for robot motion include programming, direct teaching, and teleoperation. Describing the motion in advance using a robot programming language is the simplest teaching method; however, the description becomes complicated for long-horizon tasks. However, training data can be obtained without precise modeling or parameter adjustment by having a human remotely control the robot and teaching its motion. In particular, remotely controlling the manipulator from an actual robot's perspective (the Wizard of Oz method) is desirable because it allows the robot to intuitively teach human operational skills for tasks. Workers could interact naturally with the environment by controlling the robot as if it were controlling its own body. The acquired teaching data were expected to include the information necessary for motion learning.


\textbf{Integration Learning of Sensorimotor Dynamics:}
%The learning target of our DPL concept is the time-series relationship between sensorimotor information in a system where the environment and the body dynamically interact. The training data does not have ground truth labels, and the model is trained to predict the next step robot state by taking the current robot state as input. This autoregressive learning eliminates the need for detailed design of a physical model of the environment, which is required in conventional robotics. In addition, the model can be expressed as dynamics that integrate environmental recognition and motion generation functions that span multiple modalities.
The learning target of our DPL concept is the time-series relationship between sensorimotor information in a system in which the environment and body dynamically interact. The training data did not contain ground truth labels, and the model was trained to predict the next-step robot state using the current robot state as the input. This autoregressive learning eliminates the need for a detailed design of the physical model of the environment, which is required in conventional robotics. In addition, the model can be expressed in terms of dynamics that integrate the environmental recognition and motion generation functions spanning multiple modalities.


%We construct a sensory and motor prediction system by combining multiple DNN models corresponding to the feature extraction part and the time-series prediction unit. The feature extraction unit extracts features from the sensor values acquired by the robot, and the time-series prediction unit learns sensorimotor information that integrates the extracted features and the robot state expressed by joint angles. We published the first model in 2014~\cite{Noda2014multimodal}, and since then this study has continued to expand. Each part can be connected end-to-end or learned independently. In this Section, we will clearly explain the role of each part.
We constructed a sensory and motor prediction system by combining multiple DNN models corresponding to the feature extraction and time-series prediction units. The feature extraction unit extracts features from the sensor values acquired by the robot, and the time-series prediction unit learns sensorimotor information that integrates the extracted features and the robot state expressed by joint angles. We published the first model in 2014~\cite{Noda2014multimodal}, and since then, this study has continued to expand. Each part can be connected end to end or learned independently. In this section, we explain the role of each component.


\begin{description}
\item[Feature Extraction Unit:]\mbox{}\\
%Sensor inputs that robots obtain from the real environment, such as cameras and tactile sensors, are often high-dimensional, so appropriate feature extraction is important. To appropriately learn sensorimotor information, the extracted sensor features must reflect the relationship between the robot body and the manipulated object. This is because it affects the prediction performance of motion information in the time-series learning section, such as in the case of approaching a manipulated object. Therefore, to obtain as detailed information as possible, it is desirable that the feature extractor can handle high-dimensional sensor input as is.
Sensor inputs that robots obtain from a real environment, such as cameras and tactile sensors, are often high dimensional; therefore, appropriate feature extraction is important. Sensorimotor information is appropriately learned by extracting sensor features that reflect the relationship between the robot body and the manipulated object. The extraction of sensor features affects the prediction performance of the motion information in the time-series learning section, such as when approaching a manipulated object. Therefore, to obtain as detailed information as possible, the feature extractor should handle high-dimensional sensor inputs as is.


%Here, as an easy-to-understand example, we consider image data. Raw image data contains a lot of information that is unnecessary for robot tasks, such as background, and it can become noise for training models. Many of the models we used adopted a convolutional auto-encoder (CAE) as a feature extractor. In the CAE, low-dimensional features are acquired in the middle layer by performing self-supervised learning to provide an output value that is equal to the input value. The encoded feature vector provides high-dimensional input information using fewer dimensions, making it robust against environmental noise. Furthermore, the above feature vectors can be narrowed down to information that is even more effective for task learning by introducing the attention mechanism (Section~\ref{sec_attention}).
Here, as an easy-to-understand example, we consider image data. Raw image data contain much unnecessary information for robot tasks, such as the background, and can become noise for training models. Many models adopted a convolutional autoencoder (CAE) as the feature extractor. In the CAE, low-dimensional features are acquired in the middle layer by performing self-supervised learning to provide an output value equal to the input value. The encoded feature vector provides high-dimensional input information with fewer dimensions, making it robust against environmental noise. Furthermore, the above feature vectors can be narrowed down to more effective information for task learning by introducing an attention mechanism (see Section~\ref{sec_attention}).


\item[Time-series Prediction Unit:]\mbox{}\\
%We use a Recurrent Neural Network (RNN) with a recursive connection in the middle layer as the time-series prediction unit. The RNN is a model that reflects temporal dependence, and functions as a learning-based controller that takes the current robot state as input and outputs the next robot state. The RNN takes the feature vector obtained by the sensor feature extractor and current robot state (joint angle torque value, gripper value, etc.) as input and outputs the next robot state, and learns to reduce the error between the output and correct next robot state. At this time, it is possible to construct a robust model by adjusting the ratio of how much of the model output from the previous time is used as input. It is also empirically known that predicting sensor values as a subtask contributes to the predictive performance of the model.
We used a Recurrent Neural Network (RNN) with a recursive connection in the middle layer as the time-series prediction unit. An RNN model reflects temporal dependence and functions as a learning-based controller that inputs the current robot state and outputs the next robot state. The RNN takes the feature vector obtained by the sensor feature extractor and the current robot state (joint angle torque value, gripper value, etc.) as input, outputs the next robot state, and learns to reduce the error between the output and the correct next robot state. A robust model can be constructed by adjusting the ratio of the model output to the previous time as the input. It is also empirically known that predicting sensor values as a subtask contributes to the predictive performance of a model.


%In addition to the above, the RNN can retain various dynamics obtained from the learning process of time-series data as unique internal states~\cite{Tani1998}. When the prediction error increases during robot motion generation, the RNN has the property of switching from the learned dynamics to minimize the prediction error (converge to the appropriate dynamics), and this property can be used to implement the predictive coding strategies described in the previous section. Furthermore, by representing long-horizon tasks as a combination of dynamics embedded in the RNN, it is possible to realize a composite generation of motions. The above details will be described in Section~\ref{sec_case_studies}.
Additionally, RNN can retain various dynamics obtained from the time series data learning process as unique internal states~\cite{Tani1998}. When the prediction error increases during robot motion generation, the RNN can switch from learned dynamics to minimize the prediction error (converging to the appropriate dynamics). This property can be used to implement the predictive coding strategies described in the previous section. Furthermore, the composite generation of motions can be realized by representing long-horizon tasks as a combination of dynamics embedded in the RNN. The details are presented in Section~\ref{sec_case_studies}.
\end{description}


\subsubsection{Robot Motion Generation based on Minimizing Prediction Errors}
\label{sec_generation}

%While the robot performs a task, it sequentially acquires sensor information. By performing forward computation of the model at each step, the RNN predicts the robot's state at the next time step from the input and internal context information. Then, the robot's joints are controlled using the RNN prediction as the target state. By repeating the above process online, the RNN predicts the robot's sensorimotor behavior while sequentially changing the state of each neuron in the context layer. Since our entire system is composed of light models, one of the advantages is that the calculation time and cost required for motion generation are low.
As the robot performs a task, it sequentially acquires sensor information. The RNN predicts the robot's state at the next time step by performing forward computation of the model using the input and internal context information at each step. Subsequently, the robot's joints were controlled using the RNN prediction as the target state. The RNN predicts the robot's sensorimotor behavior while sequentially changing the state of each neuron in the context layer by repeating the above process online. One of its advantages is that the calculation time and cost required for motion generation are low because our entire system is composed of light models.


%During online motion generation, the manipulated object, robot arm, sensor, and the DPL model engage each other. And, the perceptual predictions (posterior beliefs) of the model are corrected in real-time. As a result of the motion modification, the sensory input is also modified, reducing the model's prediction error. At this time, it can be observed that small trajectory corrections and operation switching are occurring through trial and error~\cite{Ichiwara2022contact,Koma2016ral}. This means transitions to other behavior attractors occur due to attraction based on prediction errors. What is important here is that trial and error is not a random process, but a correction of predicted behavior based on experience. We use the above characteristics to accomplish high-dimensional, complex, and long-horizon tasks.
During online motion generation, the manipulated object, robot arm, sensor, and DPL model engage with each other. The model's perceptual predictions (posterior beliefs) were corrected in real-time. As a result of the motion modification, the sensory input was also modified, thereby reducing the model's prediction error. It can be observed that small trajectory corrections and operation switching occur through trial and error~\cite{Ichiwara2022contact,Koma2016ral}. These observations imply that transitions to other behavior attractors occur because of attractions based on prediction errors. What is important here is that trial and error are not random processes but corrections of predicted behavior based on experience. We use these characteristics to accomplish high-dimensional, complex, and long-horizon tasks.


\subsubsection{Hierarchical Structure with Different Time-scales}
\label{sec_timescale}

%%We discuss in more detail the requirements of the RNN in the DPL concept.
%Hierarchy that considers multiple time scales of the system is also important to predictive learning. In other words, it is a system in which a model that makes predictions on a fast timescale similar to the real world and a model that makes predictions on a much slower time scale exist in parallel. We use multiple-time scale RNN (MTRNN~\cite{Yamashita2008emergence}) and LSTM~\cite{Hochreiter1997nc} in the time-series prediction unit and minimize the overall prediction error by an interaction between layers. In the lower layer, which is close to sensory input, the model predicts fast (short time constant) dynamics, and in the upper layer, the model predicts slow (long time constant) dynamics.
A hierarchy that considers multiple time scales of a system is also important for predictive learning. In other words, it is a system in which a model that makes predictions on a fast time scale similar to the real world and a model that makes predictions on a much slower time scale exist in parallel. We used a multiple-time-scale RNN (MTRNN~\cite{Yamashita2008emergence}) and LSTM~\cite{Hochreiter1997nc} in the time-series prediction unit and minimized the overall prediction error by the interaction between layers. In the lower layer, which is close to the sensory input, the model predicts fast (short time constant) dynamics, whereas in the upper layer, the model predicts slow (long time constant) dynamics.


%In the above model, sensory prediction errors are absorbed by changes in the state of the lower layer. At this time, the upper layer generates predictions based on long-term intentions and provides the lower layer with constraints to maintain this state. When a sensory prediction error can no longer be absorbed by the lower layer, the error affects the upper layer, causing a state change (change of intention) in the upper layer. This structure in which the upper and lower layers interact with each other enables consistent yet adaptive behavior generation~\cite{Suzuki2018ral,Saito2021}.
In the above model, the sensory prediction errors were absorbed by changes in the state of the lower layer. The upper layer generates predictions based on long-term intentions, giving the lower layer constraints to maintain this state. When the lower layer can no longer absorb a sensory prediction error, it affects the upper layer, causing a state change (change in intention) in the upper layer. This structure, in which the upper and lower layers interact, enables consistent yet adaptive behavior generation~\cite{Suzuki2018ral,Saito2021}.


\subsubsection{Bottom-up and Top-down Attention}
\label{sec_attention}

%Next, we describe the attention mechanism involved in modifying prediction accuracy. In general, attention mechanisms efficiently allocate a model's resources by concentrating the processing area on a specific target. In robot tasks, it is possible to provide robustness against sudden changes in lighting conditions, placement of multiple objects, and movement of target objects that cause noise. Usually, attention in image processing is generated bottom-up directly from the current input image~\cite{Vaswani2017nips}. For example, methods include emphasizing features such as brightness, color, direction, and orientation, or generating noteworthy keys, queries, and values in self-attention.
Next, we describe the attention mechanism involved in modifying prediction accuracy. Generally, attention mechanisms efficiently allocate a model's resources by concentrating the processing area on a specific target. In robotic tasks, it is possible to provide robustness against sudden changes in lighting conditions, placement of multiple objects, and movement of target objects that cause noise. Typically, attention in image processing is generated bottom-up directly from the current input images~\cite{Vaswani2017nips}. For example, methods emphasize brightness, color, direction, and orientation or generate noteworthy keys, queries, and values in self-attention.


%In the DPL framework, we are conscious of guiding the attention area by prediction. By predicting multiple attention points within an image by using spatial softmax~\cite{finn2016deep}, we limit the areas that can be attended to in the image and have no choice but to ignore other areas. We use the attention point and sensorimotor prediction error as a loss function, and as a result, the attention point is self-organized in that it changes depending on one's actions. In addition to the above bottom-up attention, top-down attention generates attention areas based on the intentions of the upper layers. In \cite{Ito2022integrated}, it is possible to confirm the background other than the object being attended to disappears in a predicted image generated as a posteriori belief from a verbal instruction. Please refer to our related work~\cite{Ito2022integrated,Ichiwara2022contact,Ichiwara2023ral} for details on this attention mechanism.
In the DPL framework, we are conscious of guiding the attention areas through predictions. By predicting multiple attention points within an image using spatial softmax~\cite{finn2016deep}, we limit the areas that can be attended to in the image and have no choice but to ignore others. We use the attention point and sensorimotor prediction error as a loss function. As a result, the attention point is self-organized, changing depending on one's actions. In addition to bottom-up attention, top-down attention generates attention areas based on the intentions of the upper layer. In \cite{Ito2022integrated}, it was possible to confirm that the background, other than the object being attended to, disappeared in a predicted image generated as a posteriori belief from verbal instruction. Please refer to our related work~\cite{Ito2022integrated,Ichiwara2022contact,Ichiwara2023ral} for details regarding this attention mechanism.



\subsection{Discussion of Proposed Concept}
\subsubsection{Perspectives on Embodiment and Dynamic Cognitive Systems}

%In the fields based on physicality, such as behavior-based robotics~\cite{Brooks1991intelligence}, it is believed that intelligent movements emerge from the coupling of the environment, body, and nervous system. In the above fields, the importance of the interaction between the body and the environment is particularly emphasized. For example, in a passive walking mechanism, walking motion is generated through the interaction between a well-designed bipedal link model and a slope~\cite{Garcia1996}. Considering the body's flexibility, mechanisms, and sensor placement, the body is thought to potentially reserve a variety of actions. At this time, it is said that the nervous system only needs to perform the minimum functions, and in the example of walking, it is sufficient for the oscillator system to function.
In fields based on physicality, such as behavior-based robotics~\cite{Brooks1991intelligence}, intelligent movements are believed to emerge from the coupling of the environment, body, and nervous system. The above fields emphasize the importance of the interaction between the body and the environment. For example, in a passive walking mechanism, the walking motion is generated through the interaction between a well-designed bipedal link model and a slope~\cite{Garcia1996}. Considering the body's flexibility, mechanisms, and sensor placement, it is believed to reserve a variety of actions. At this time, the nervous system only needs to perform minimal functions, and in the example of walking, it is sufficient for the oscillator system to function.


%With the above coupling, it is also possible to provide various states of attraction (attractors) to the nervous system. The DPL concept can be seen as a framework in which the above nervous system part is constructed using a deep learning model. For example, there is a discussion of dynamic cognitive systems based on the perspective of dynamical systems proposed by Tani~\cite{Tani1996model}. Here, the brain is not considered simply as an input/output system, but as a predictive system for the world. During interactions through the body, the brain sometimes forms a specific attractor structure with the world, and when the discrepancy becomes large, this connection breaks. After that, it transitions to a new attractor, sometimes through a phase transition, sometimes through a chaotic state. As this process is repeated, cognition changes dynamically without converging to a fixed state. Tani et al. have realized this process using RNN and a real robot system, and have been conducting interesting research for many years in an attempt to understand the brain cognitive process~\cite{tani2016exploring}. Although the above study was proposed before the discussion of predictive coding and the free energy principle, it has many commonalities with our DPL concept.
This coupling can provide various states of attraction (attractors) to the nervous system. The DPL concept can be considered a framework in which the above nervous system part is constructed using a deep learning model. For example, there is a discussion of dynamic cognitive systems based on the perspective of dynamical systems proposed by Tani~\cite{Tani1996model}. Here, the brain is not simply considered an input/output system but a predictive system for the world. During interactions with the body, the brain sometimes forms a specific attractor structure with the world; this connection breaks when this discrepancy becomes large. Subsequently, it transitions to a new attractor, sometimes through a phase transition or a chaotic state. As this process is repeated, cognition changes dynamically without converging to a fixed state. Tani et al. realized this process using RNN and a real robot system and have been conducting interesting research for many years, attempting to understand the brain cognitive process~\cite{tani2016exploring}. Although the above study was proposed before the discussion on predictive coding and the free-energy principle, it has many similarities with our DPL concept.


\subsubsection{Contrasting with Reinforcement Learning}

%Finally, we discuss the main differences between reinforcement learning and predictive coding. The discussion here is just a comparison with simple reinforcement learning. As mentioned in the previous chapter, one of the distinctive charms of reinforcement learning is that learning goals can be set freely (for example, scalar values, etc.). The main objective of predictive coding is to minimize the prediction error when the model connects the body and the world. Therefore, it is difficult in principle to set other learning goals in predictive coding. Therefore, unless special measures are taken, it is difficult in principle to set other learning goals in predictive coding.
Finally, we discuss the main differences between reinforcement learning and predictive coding. The discussion here is a simple comparison with simple reinforcement learning. As mentioned in the previous section, one of the distinctive characteristics of reinforcement learning is that learning goals can be freely set (for example, scalar values, etc.). The main objective of predictive coding is to minimize prediction errors when the model connects the body and the world. Therefore, setting other learning goals for predictive coding is difficult. Therefore, setting other learning goals for predictive coding will be difficult unless special measures are taken.


%In model-based reinforcement learning, a method has been proposed to learn a "world model" by repeating interactions between actions and states~\cite{Ha2018world}. The method learns actions and visual field images as a sequence of abstracted latent spaces. This is similar in concept to predictive coding in that it targets the world as a prediction target. After learning has progressed sufficiently, the optimal behavior based on the world model can be obtained by searching through trial and error for the behavioral policy that yields the maximum reward from the model.
In model-based reinforcement learning, a method was proposed for learning a real-world model" by repeating the interactions between actions and states~\cite{Ha2018world}. This method learns actions and visual field images as a sequence of abstracted latent spaces. The learning process is similar to predictive coding in that it targets the world as a prediction target. After learning has progressed sufficiently, the optimal behavior based on the world model can be obtained by searching through trial and error for a behavioral policy that yields the maximum reward from the model.


%On the other hand, predictive coding assumes that the world is open and that the environment model learning is constantly progressing. While learning is insufficient, actions are generated that are expected to reduce future prediction errors (expected free energy). By repeating the experience of reaching the goal state, the sequence of actions becomes a habit as a prior belief. Thus, actions with high prediction accuracy are selected. In this process of habituation, the generated action sequence (virtual expectations) can be thought of as trial and error. In other words, the behavior learned in this framework is not necessarily optimal for a given objective index. The search that accompanies learning has the aspect of not only optimization but also diversification of possible adaptations. The paradox of trying to adapt to the environment but not being able to do so completely leads to the exploration of various actions.
In contrast, predictive coding assumes that the world is open and that environmental model learning is constantly progressing. When the learning is insufficient, actions expected to reduce future prediction errors (expected free energy) are generated. By repeating the experience of reaching the goal state, the sequence of actions becomes a habit, as a prior belief. Thus, actions with high prediction accuracy were selected. The generated action sequence (virtual expectations) can be considered trial and error in this habituation process. In other words, the behavior learned in this framework is not necessarily optimal for a given objective index. The search that accompanies learning involves optimization and diversification of possible adaptations. The paradox of trying to adapt to the environment but not being able to do so leads to the exploration of various actions.





\section{Case Studies}

%In this section, we summarize the published studies that we have developed using the DPL concept from two perspectives: scalability by combining multiple motions, and application to complex manipulation tasks.
This section summarizes the published studies we developed using the DPL concept from two perspectives: scalability by combining multiple motions and application to complex manipulation tasks.



\subsection{Compound Generation of Robot Motions}
\label{sec_case_studies}

%As mentioned in the first section, approaches based on large-scale data and trials have high training costs, so it is necessary to perform a wide range of tasks with limited robot motion data. Here, since many practical tasks are essentially long-horizon tasks that consist of multiple steps, an approach that expresses them by dividing and combining short tasks (subtasks) is effective~\cite{Sacerdoti1974ai,Wolpert1998,Schaal2003isrr}. This kind of "compound generation of robot motions" has the following advantages.
As mentioned in the first section, approaches based on large-scale data and trials have high training costs; therefore, it is necessary to perform a wide range of tasks using limited robot motion data. It is an effective approach to express practical tasks by dividing and combining short tasks (subtasks)because many are essentially long-horizon tasks consisting of multiple steps~\cite{Sacerdoti1974ai,Wolpert1998,Schaal2003isrr}. This type of "compound generation of robot motions" has the following advantages:
\begin{itemize}
    \item Using subtasks included in a long-term task for another task (module reuse)
    \item Executing new tasks by combining subtasks (extending functionality)
\end{itemize}
%With the above background, this section discusses the elements required to combine multiple motions and introduces robot applications with the DPL concept.
With this background, this section discusses the elements required to combine multiple motions and introduces robot applications using the DPL concept.


\subsubsection{Multi-Module Switching based on Prediction Error}

%One of the simple approaches to compound-generating robot motions is the cooperation of modularized functions~\cite{Haruno2001,Wolpert1998}. Generally, a robot system consists of a group of multiple modules, and the output is determined by combining them. Even in robot learning, it can be assumed that by combining or switching models with different functions in parallel, the overall task performance can be expanded without increasing the computational load. Furthermore, by having a distributed system configuration, it is possible to continuously add functionality by adding modules.
A simple approach to compound-generating robot motion involves the cooperation of modularized functions~\cite{Haruno2001,Wolpert1998}. Generally, a robotic system consists of multiple modules, and the output is determined by combining them. Even in robot learning, it can be assumed that by combining or switching models with different functions in parallel, the overall task performance can be improved without increasing the computational load. Furthermore, with a distributed system configuration, it is possible to continuously add functionality by adding modules.


%One method to achieve the above in robot learning is to switch subtask modules based on the prediction error of the model. As mentioned in the previous Section, RNNs can retain various dynamics obtained from the learning process as unique internal states. When generating robot motion, RNN predicts from the internal state and current sensorimotor information, and, prediction errors tend to increase in untrained states. Therefore, when the prediction error increases during motion generation, the RNN has the property of switching from the dynamics to minimize the prediction error, and this can be used to switch modules. When a large prediction error that cannot be handled in real-time occurs, it cannot be handled by the state transition of the attracter embedded in individual RNNs. In that case, it is necessary to adjust the overall sensory prediction and motion generation by coordinating multiple modules according to the amount of the prediction error.
One method for achieving this in robot learning is to switch the subtask modules based on the model's prediction error. As mentioned in the previous Section, RNNs can retain the various dynamics obtained from the learning process as unique internal states. When generating robot motion, RNN predicts from the internal state and current sensorimotor information, and the prediction errors tend to increase in untrained states. Therefore, when the prediction error increases during motion generation, the RNN switches from dynamics to minimize the prediction error, which can be used to switch modules. When a large prediction error that cannot be handled in real-time occurs, it cannot be handled by the state transition of the attracter embedded in the individual RNNs. In this case, it is necessary to adjust the overall sensory prediction and motion generation by coordinating multiple modules according to the prediction error.


%\cite{Ito2022sr} proposed a method that takes actions to reduce prediction errors in door-passing tasks by introducing multiple DPL models and a new mechanism that switches them according to prediction errors (Figure~\ref{fig_doorrobot}). The model was trained for each subtask, such as ``opening a door'' or ``passing through,'' and then combined these models to perform a series of tasks. Each model predicted the near future camera image from sensor information. By comparing it with the real camera image, the method calculated the change in confidence value, which indicates how accurately the robot can work in the real environment, over time. Furthermore, the robot executed actions appropriate to the situation by autonomously selecting the model with the highest confidence. Since the robot performed these calculations in real-time, it could handle complex tasks, without having to accurately design the timing of switching motions or the flow of motion. As a result of real robot experiments, we confirmed that it is possible to generate tasks in response to unlearned door positions and changes in door panel color and pattern\footnote{https://youtu.be/pdOA-PeYO9Q}.
\cite{Ito2022sr} proposed a method that takes actions to reduce prediction errors in door-passing tasks by introducing multiple DPL models and a new mechanism that switches them according to prediction errors (Figure~\ref{fig_doorrobot}). The model was trained for each subtask, such as ``opening a door'' or ``passing through,'' and then combined these models to perform a series of tasks. Each model predicts a near-future camera image from the sensor information. The method calculates the change in the confidence value, which indicates how accurately the robot can work in a real environment over time by comparing it with a real camera image. Furthermore, the robot executes actions that are appropriate for the situation by autonomously selecting the model with the highest confidence. It could handle complex tasks without having to accurately design the timing of the switching motions or the flow of motion because the robot performed these calculations in real time. Based on real robot experiments, we confirmed that generating tasks in response to unlearned door positions and changes in the door panel color and pattern\footnote{https://youtu.be/pdOA-PeYO9Q} is possible.


%-------------------- Fig --------------------
\begin{figure}[t]
\centering
\includegraphics[width=73mm]{DoorRobotv2.eps}
\caption{
Real-time switching of multiple motion generation models based on prediction errors~\cite{Ito2022sr}.
}
\label{fig_doorrobot}
\hspace{10mm}
\end{figure}


%On the other hand, \cite{Suzuki2021ral} worked on responding to robot states that are not included in the training data by utilizing multiple control modules with different characteristics. This is an attempt to deal with the extrapolation of motion datasets (undefined behaviors) that cannot be predicted by a learning-based controller. In the proposed method, we prepared a model-based controller that compensates the output of the RNN. And, we appropriately switched models according to the error between the predicted trajectory and the actual trajectory. Assuming that the RNN embedded the most recent information necessary for task execution in its internal state, we trained additional tasks to predict past trajectories from the internal state. By adding a new neuron that predicts additional tasks to the middle layer of the RNN, it is possible to add error recovery functionality without affecting the performance of the main task learned. In real robot experiments, we confirmed that by combining this method with the dynamic design of the RNN internal state described in the next section, it is possible to perform error recovery in response to multiple types of disturbances (environment contact and changes) from collaborative workers.
Contrarily, \cite{Suzuki2021ral} worked on responding to robot states not included in the training data by utilizing multiple control modules with different characteristics. This approach attempted to extrapolate motion datasets (undefined behaviors) that cannot be predicted using a learning-based controller. In the proposed method, we prepared a model-based controller that compensated for the output of the RNN. In addition, the models were appropriately switched based on the error between the predicted and actual trajectories. Assuming that the RNN embeds the most recent information necessary for task execution in its internal state, we train additional tasks to predict past trajectories from the internal state. By adding a new neuron that predicts additional tasks to the middle layer of the RNN, an error recovery functionality can be added without affecting the performance of the main learned task. In real robot experiments, we confirmed that, by combining this method with the dynamic design of the RNN internal state described in the next section, it is possible to recover errors in response to multiple types of disturbances (environmental contact and changes) from collaborative workers.


%In general robot systems, it is necessary to separately program responses to unexpected situations, which requires a lot of development cost and adjustment work. Additionally, when an unexpected disturbance occurs in the work environment, the robot has to take time to recognize the situation and replan, which causes the robot to move slowly. Our proposed methods~\cite{Ito2022sr,Suzuki2021ral} can respond to sudden situational changes and changes in operating procedures by switching multiple controllers according to prediction errors, and are expected to solve the above problems.
In general robot systems, it is necessary to separately program responses to unexpected situations, which require significant development costs and adjustment work. In addition, when an unexpected disturbance occurs in the work environment, the robot must take time to recognize the situation and replan, which causes it to move slowly. Our proposed methods~\cite{Ito2022sr,Suzuki2021ral} can respond to sudden situational changes and changes in operating procedures by switching multiple controllers according to prediction errors and are expected to solve the above problems.


\subsubsection{Switching of Multiple Behaviours by Designing the Dynamical Systems}

%Next, we discuss the dynamical system of the RNN as a method for handling local composite generation within a single module. The method described in the previous section focused on switching between modules, however, considering scalability, it is desirable to have a mechanism that can also switch operations within individual modules. This is because when a task is further subdivided, each module is required to handle multiple subtasks~\cite{Nair2019iclr,Yu2019iros}. Previous studies~\cite{Yin2021iclr,Garrett2021cras} have also investigated motion planning that combines subtasks. To connect multiple robot motions, it is necessary to express their dynamics in a combinable format within a module. By selecting appropriate dynamics of the RNN according to the sensory input, it becomes possible to execute long-term task operations in response to environmental changes. We implemented the above by explicitly embedding attractor dynamics in subtask units into the RNN.
Next, we discuss the dynamic system of the RNN as a method for handling local composite generation within a single module. The method described in the previous section focuses on switching between modules; however, considering scalability, it is desirable to have a mechanism that can switch operations within individual modules. This need arises because when a task is subdivided, each module must handle multiple subtasks~\cite{Nair2019iclr,Yu2019iros}. Previous studies~\cite{Yin2021iclr,Garrett2021cras} also investigated motion planning that combines subtasks. Expressing their dynamics in a combined format within a module is necessary to connect multiple robot motions. Long-term task operations can be executed in response to environmental changes by selecting the appropriate dynamics of the RNN according to the sensory input. We implemented the above by explicitly embedding the attractor dynamics in the subtask units into an RNN.


%In \cite{Kase2018icra}, to combine multiple subtasks, we intentionally embedded an attractor structure inside the model by using a loss function in the internal state of the RNN (Figure~\ref{fig_kase_attractor}). By adding a constraint loss to have a shared state that can be switched between the attractors of each subtask, the attractor switches its dynamics to reduce the prediction error. In the experiment, by combining multiple subtasks for a Put-in-Box task, we achieved motion generation in an arbitrary order of subtasks\footnote{https://youtu.be/QwogmDqOxfE}. Furthermore, we expanded the loss function of the internal state of RNN and designed an arbitrary attractor structure to increase the degree of freedom in switching postures~\cite{Kase2019iros}.
In \cite{Kase2018icra}, to combine multiple subtasks, we intentionally embedded an attractor structure inside the model using a loss function in the internal state of the RNN (Figure~\ref{fig_kase_attractor}). The attractor switches its dynamics to reduce the prediction error by adding a constraint loss to obtain a shared state that can be switched between the attractors of each subtask. In this experiment, combining multiple subtasks for a put-in-box task generated motions in an arbitrary order of subtasks\footnote{https://youtu.be/QwogmDqOxfE}. Furthermore, we expanded the loss function of the internal state of RNN and designed an arbitrary attractor structure to increase the degrees of freedom in switching postures~\cite{Kase2019iros}.


%-------------------- Fig --------------------
\setlength\textfloatsep{5pt}
\begin{figure}[t]
\centering
\includegraphics[width=\columnwidth]{icra2018_attractor_v2.eps}
\caption{
Embedding attractor dynamics into the internal state of RNN~\cite{Kase2018icra}.
}
\label{fig_kase_attractor}
\hspace{10mm}
\end{figure}


%\cite{Suzuki2018ral} introduced a hierarchical structure of task instruction input and RNN to learn how to switch subtasks during a long-horizon task. By inputting a task instruction vector at the designed point attractor, it is possible to respond to motion switching in ambiguous situations that cannot be determined based on sensor information alone. In the experiment, we verified a clothes-folding task that consisted of multiple subtasks and confirmed that the robot could perform according to the instructions. Furthermore, in the internal state of the RNN, instructions were embedded in the long-term time scale layer, and motion was embedded in the short-term time scale layer as an attractor structure. This indicated that the task context, language, and image input were appropriately structured to combine.
\cite{Suzuki2018ral} introduced a hierarchical structure of task instruction input and RNN to learn how to switch subtasks during a long-horizon task. Inputting a task instruction vector to the designed point attractor enables the response to motion switching in ambiguous situations that cannot be determined based on sensor information alone. In the experiment, we verified a clothes-folding task that consisted of multiple subtasks and confirmed that the robot could perform it according to the instructions. Furthermore, in the internal state of the RNN, instructions were embedded in the long-term timescale layer, and motion was embedded in the short-term timescale layer as an attractor structure. These embeddings indicated that the task context, language, and image inputs were appropriately combined.


%Although these motion combination/switching methods using attractor structures require the design of switching postures for subtasks, they can incorporate the human's intentions to some extent into the RNN output. Combined with the grounding with the language described below, it is expected that the robot will be able to switch its motions more flexibly.
Although these motion combination/switching methods using attractor structures require the design of switching postures for subtasks, they can incorporate human intentions to some extent into the RNN output. The robot is expected to be able to switch its motions more flexibly, combined with grounding, using the language described below.



\subsubsection{Language-Conditioned Motion Generation}

%To express more primitive combinations of actions, it is necessary to ground each segmented subtask and language expression. By learning that the given language instructions and individual tasks are semantically grounded, the RNN expresses the internal dynamics with the semantic and syntactic structure of the language system as a type of constraint. This internal representation can be extended to use cases where actions are sequentially generated while communicating with collaborators. Based on the above idea, we aimed to inductively express the relationship between motions and symbols contained in sensorimotor information~\cite{Harned1990pd} within a model.
Grounding each segmented subtask and language expression is necessary to express more primitive action combinations. The RNN expresses internal dynamics with the semantic and syntactic structure of the language system as a constraint by learning that the given language instructions and individual tasks are semantically grounded. This internal representation can be extended to cases where actions are sequentially generated while communicating with collaborators. Based on the above idea, we aim to inductively express the relationship between motions and symbols contained in sensorimotor information~\cite{Harned1990pd} within a model.



\textbf{Grounding of Language and Robot Motion:}
%In general, the correspondence between language, which is a discrete expression, and robot motions in a continuous space changes depending on the past context and the situation on the spot. For example, even with a simple instruction such as "Hit the blue bell," the action to accomplish this will vary depending on where the blue bell is placed, what posture the robot is in, etc. Robots need to acquire internal representations that can be used in a generalized manner even in unknown situations, so that they can achieve such ambiguous and one-time grounding each time.
In general, the correspondence between language, which is a discrete expression, and robot motion in a continuous space changes depending on the past context and situation on the spot. For example, even with simple instructions such as "Hit the bluebell," the action to accomplish this will vary depending on where the bluebell is placed, what position the robot is in, etc. Robots must acquire internal representations that can be generalized, even in unknown situations, to achieve such ambiguous and one-time grounding each time.


%Our DPL framework makes it possible to embed the understanding of language instructions and the motions into a set of attractor structures. \cite{Yamada2016fn,Yamada2017fn} acquire cyclic attractors that express the above episodic repetition as the internal dynamics by having RNN predict and learn a sequence that repeats multiple episodes consisting of waiting, verbal instructions, and robot motion generation. (Figure~\ref{fig_language_reduce}). When a verbal instruction and a camera image are input to the RNN at a fixed point attractor, the internal state of the RNN changes in a branching manner according to the local situation. After the language input is completed, robot motions are generated according to the meaning of the instructions in the environment. Then, once the motion is completed, the internal state converges to a fixed point again and waits for the next instruction. We demonstrated the generalization ability to generate appropriate actions for new situations without having to teach every combination of environmental conditions and verbal instructions.
Our DPL framework enables embedding an understanding of language instruction and motion into a set of attractor structures. \cite{Yamada2016fn,Yamada2017fn} acquire cyclic attractors that express the above episodic repetition as the internal dynamics by having RNN predict and learn a sequence that repeats multiple episodes of waiting, verbal instructions, and robot motion generation. (see Figure~\ref{fig_language_reduce}). When a verbal instruction and a camera image are input to the RNN at a fixed-point attractor, the internal state of the RNN changes in a branching manner according to the local situation. After the language input was completed, robot motions were generated according to the meaning of the instructions in the environment. Once the motion was complete, the internal state converged to a fixed point and waited for the next instruction. We demonstrated the generalization ability to generate appropriate actions for new situations without having to teach every combination of environmental conditions and verbal instructions.


%-------------------- Fig --------------------
\setlength\textfloatsep{5pt}
\begin{figure}[t]
\centering
\includegraphics[width=\columnwidth]{rsj2022_language_reduce.eps}
\caption{
Grounding of language instructions and motion generation in the designed internal dynamics of RNN~\cite{Yamada2016fn}.
}
\label{fig_language_reduce}
\hspace{10mm}
\end{figure}


%Furthermore, it is also possible to extend this model to a bidirectional conversion model that includes conversion from task instructions to actions and from actions to linguistic explanations. In \cite{Yamada2018ral}, we prepared two autoencoders, one for language and one for sensorimotor sequences, and performed learning by setting constraints such that the feature vectors of the corresponding sentences and motions were close to each other. After the training phase, the model achieved two-way translation from instructions to robot motions and robot motions to linguistic explanations via feature vectors. \cite{Toyoda2021ral,Toyoda2022ral} added that model a word distributed representation conversion module that utilizes the Word2vec model, allowing it to perform appropriate operations even for instructions that include new words similar to trained words.
Furthermore, extending this model to a bidirectional conversion model that includes conversion from task instructions to actions and from actions to linguistic explanations is possible. In \cite{Yamada2018ral}, we prepared two autoencoders, one for language and one for sensorimotor sequences. We performed learning by setting constraints such that the feature vectors of the corresponding sentences and motions were close to each other. After the training phase, the model achieved a two-way translation from instructions to robot motions and robot motions to linguistic explanations via feature vectors. \cite{Toyoda2021ral,Toyoda2022ral} added that model a word distributed representation conversion module that utilizes the Word2vec model, allowing it to perform appropriate operations even for instructions that include new words similar to trained words.


\textbf{Integration with Task Planning:}
%In the above studies, we introduced methods to give instructions in the form of language and integrate the robot's motions. Linguistic instructions are given word by word, but general-purpose planning description languages (e.g. Planning Domain Description Language: PDDL), which are used in conventional automatic planning methods, can be used as task instructions. \cite{Kase2020icra} proposed a learning method to connect descriptive language that defines the environment and images (Figure~\ref{fig_kase2022icra}). It designed motions using any automatic planning method based on discrete symbolic representations extracted from images. Although conventional PDDL makes it difficult to respond to changes in the environment because it is planned, our method makes it possible to dynamically determine the robot's actions.
In the above studies, we introduced methods to provide instructions in the language form and integrate the robot's motions. Linguistic instructions are given word by word; however, general-purpose planning description languages (e.g., planning domain description language: PDDL) used in conventional automatic planning methods can be used as task instructions. \cite{Kase2020icra} proposed a learning method to connect descriptive language that defines the environment and images (Figure~\ref{fig_kase2022icra}). They designed motions using an automatic planning method based on discrete symbolic representations extracted from images. Although conventional PDDL makes it difficult to respond to planned environmental changes, our method makes it possible to determine the robot's actions dynamically.


%In recent years, with the development of large language models (LLM~\cite{Zhao2023}), they have been actively applied to robot motion generation~\cite{Shridhar2021a,Liang2022a,Brohan2022a,Driess2023a}. Especially in task planning, it is possible to convert linguistic instruction to multiple subtask instructions in table format~\cite{Ahn2022,huang2022language} and code format~\cite{Vemprala2023,codeaspolicies2022}. Although our studies introduced in this section deal with limited language formats, we plan to expand them by incorporating LLM technology~\cite{Hori2024}.
In recent years, large language models (LLM~\cite{Zhao2023}) have been actively applied in robot motion generation~\cite{Shridhar2021a,Liang2022a,Brohan2022a,Driess2023a}. Especially in task planning, it is possible to convert linguistic instructions into multiple subtask instructions in table format~\cite{Ahn2022,huang2022language} and code format~\cite{Vemprala2023,codeaspolicies2022}. Although the studies introduced in this section deal with limited language formats, we plan to expand them by incorporating LLM technology~\cite{Hori2024}.


%-------------------- Fig --------------------
\begin{figure}[t]
\centering
\includegraphics[width=\columnwidth]{icra2020_overview.eps}
\caption{
Grounding with a generic planning language and extracted image features~\cite{Kase2020icra}.
}
\label{fig_kase2022icra}
\hspace{10mm}
\end{figure}    



\subsection{Application to Various Robot Tasks}
\subsubsection{Deformable Object Manipulation}

%We also verified the versatility of the DPL concept by applying it to difficult manipulation tasks. A typical example is the task of manipulating deformable objects. Manipulating with flexible objects such as cloth and strings that change shape upon contact makes it difficult to describe the entire situation in a program. \cite{Koma2016ral} generalized the motion of folding a towel by performing integrated learning based on image-motion information acquired through the robot's work experience. As a result, the robot could repeatedly perform tasks that correspond to object position shifts and different objects (such as books) from limited training data. Unlike conventional methods, our method does not require high-cost processing such as modeling or image processing, and it is possible to generate motion at a natural speed\footnote{https://youtu.be/YH1TrL1q6Po}. Our concept does not depend on tasks or robot hardware, so by implementing each function as a component, we easily reuse the implemented system in response to changes in requirements~\cite{Kanamura2021sensors}. In previous studies, we realized  the in-air knotting task~\cite{Suzuki2021iros} and the buttoning task~\cite{Fujii2022sii} by using the same model.
We verified the versatility of the DPL concept by applying it to difficult tasks. A typical example is the manipulation of deformable objects. Manipulating flexible objects such as cloths and strings that change shape upon contact makes it difficult to describe the entire situation in a program. \cite{Koma2016ral} generalized the motion of folding a towel by performing integrated learning based on image-motion information acquired through the robot's work experience. As a result, the robot could repeatedly perform tasks corresponding to object position shifts and different objects (such as books) using limited training data. Unlike conventional methods, our method does not require high-cost processing such as modeling or image processing, and it is possible to generate motion at a natural speed\footnote{https://youtu.be/YH1TrL1q6Po}. Our concept does not depend on tasks or robot hardware; therefore, by implementing each function as a component, we can easily reuse the implemented system in response to changes in requirements~\cite{Kanamura2021sensors}. We realize the in-air knotting task~\cite{Suzuki2021iros} and the buttoning task~\cite{Fujii2022sii} using the same model in previous studies.


%It has been confirmed that including detailed object information such as proximity sensors and tactile sensors as input to the model improves prediction performance. \cite{Ichiwara2022contact} introduced an attention mechanism that extracted information (such as the position and direction of the object) from the robot's visual and tactile information (Figure~\ref{fig_attention_bag}). In this task, it is necessary to generate an appropriate trajectory in real time according to the state of the bag and robot hand. The game starts with grabbing randomly placed zippers, so the gripping position and posture constantly change. The number of training sequences was 36, which is extremely small compared to general learning methods. In the experiment, it was confirmed that the perceptual predictions (predicted images, tactile states, and motions) of the model were corrected in real time by the attention mechanism.
It has been confirmed that including detailed object information, such as proximity and tactile sensors, as input to the model improves the prediction performance. \cite{Ichiwara2022contact} introduced an attention mechanism that extracted information (such as the position and direction of the object) from the robot's visual and tactile information (Figure~\ref{fig_attention_bag}). In this task, generating an appropriate trajectory in real-time is necessary according to the state of the bag and robot hand. The game starts by grabbing randomly placed zippers so that the gripping position and posture constantly change. The number of training sequences was 36, which is extremely small compared with general learning methods. The experiment confirmed that the model's perceptual predictions (predicted images, tactile states, and motions) were corrected in real time using an attention mechanism.


%The sense of touch also plays an important role in the emergence of diversity generated from exploration based on intrinsic motivation. Predictability and operability are criteria of the learning, however, they do not indicate the possibility of further exploration. We also proposed searching for new actions based on tactile perceptual information~\cite{Mori2020} and active inference based on the internal characteristics of the manipulated object~\cite{Saito2021,Saito2023a}. These studies addressed the next important issue to be solved in our DPL concept from the perspective of the emergence of various movements based on exploration.
Touch also plays an important role in the emergence of diversity generated from exploration based on intrinsic motivation. Predictability and operability are the criteria for learning; however, they do not indicate the possibility of further exploration. We also proposed searching for new actions based on tactile perceptual information~\cite{Mori2020} and active inference based on the internal characteristics of the manipulated object~\cite{Saito2021,Saito2023a}. These studies address the next important issue to be solved in our DPL concept from the perspective of the emergence of various movements based on exploration.


%-------------------- Fig --------------------
\setlength\textfloatsep{5pt}
\begin{figure}[t]
\centering
\includegraphics[width=\columnwidth]{attention_bag.eps}
\caption{
Generated bag opening task using attention mechanism with vision and tactile sensors~\cite{Ichiwara2022contact}.
}
\label{fig_attention_bag}
\hspace{10mm}
\end{figure}   



\subsubsection{Neuro-Robotics Experiments}

%We also studied to use a similar hierarchical predictive model as a robot simulation for neurodevelopmental disorders. We reproduced the functional disconnection between layers~\cite{idei2021paradoxical} and the imbalance of excitability and inhibition in neurons~\cite{idei2020homogeneous}. The results of these studies show that when a real robot causes an abnormality in sensory accuracy or even upper-layer prediction accuracy, this is observed as a behavioral abnormality. Although this paper does not go into detail, the above studies are significant in that they extend the DPL concept from a perspective not limited to task learning.
We also used a similar hierarchical predictive model for robot simulation of neurodevelopmental disorders. We reproduced the functional disconnection between layers~\cite{idei2021paradoxical} and the imbalance between excitability and inhibition in neurons~\cite{idei2020homogeneous}. These studies show that when a real robot causes an abnormality in sensory accuracy or even upper-layer prediction accuracy, it is observed as a behavioral abnormality. Although this paper does not go into detail, the studies mentioned above are significant because they extend the DPL concept from a perspective not limited to task learning.


\section{Conclusion}

%In this paper, we introduced deep predictive learning which is the robot motion learning concept. We also summarized the background technology of predictive coding, related studies, and robot applications of our proposed concept. The three robot application examples described in terms of the composition of multiple motions are as follows; switching between multiple modules based on prediction errors, designing the dynamical systems inside the DNN model, and acquiring an integrated language--motion representation. In addition, we have shown their applicability to a wide variety of robotic tasks, including flexible object manipulation. Although each method has been developed separately, there are many common technical elements and it is expected that they will be integrated in the future. However, there are challenges to improving accuracy and applying the system to even more difficult tasks. We plan to revisit the proposed concept from the perspective of cognitive robotics, which is the origin of the concept, and apply the method to longer-term and more diverse tasks in line with real-world problems.
In this study, we introduce deep predictive learning, a concept of robot motion learning. We also summarize the background technology of predictive coding, related studies, and robotic applications of the proposed concept. The three robot application examples described regarding the composition of multiple motions are as follows: switching between multiple modules based on prediction errors, designing dynamic systems inside the DNN model, and acquiring an integrated language-motion representation. In addition, we demonstrate their applicability to various robotic tasks, including flexible object manipulation. Although each method has been developed separately, many common technical elements are expected to be integrated in the future. However, there are challenges in improving the accuracy and applying the system to more difficult tasks. We plan to revisit the proposed concept from the perspective of cognitive robotics the concept's origin, and apply the method to long-term and more diverse tasks in line with real-world problems.



\section*{ACKNOWLEDGMENT}
The authors would like to thank Dr. Hayato Idei for his comments on the description of predictive coding in this manuscript.
This work was supported by JST Moonshot R\&D Grant Number JPMJMS2031, JST ACT-X GrantNumber JPMJAX190I, Hitachi Ltd. Japan.

\bibliographystyle{ieeetr}
\bibliography{refs2}

\end{document}

